% Part: first-order-logic 
% Chapter: axiomatic-deduction 
% Section: deduction-theorem-quantifiers

\documentclass[../../../include/open-logic-section]{subfiles}

\begin{document}

\olfileid[pt]{fol}{axd}{ddq}

\olsection{O teorema da dedução com quantificadores}

\begin{thm}[Teorema da dedução]
\ollabel{thm:deduction-thm-q} Se $\Gamma \cup \{!A\} \Proves !B$,
então $\Gamma \Proves !A \lif !B$.
\end{thm}

\begin{proof}
Procedemos novamente por indução sobre o comprimento
da !!{derivation} de $!B$ a partir de $\Gamma \cup \{!A\}$.

A demonstração da base da indução é idêntica à apresentada na
demonstração de~\olref[ded]{thm:deduction-thm}.

Para o passo indutivo, suponha novamente que a !!{derivation} de $!B$
a partir de $\Gamma \cup \{!A\}$ termina com um passo~$!B$ justificado por
uma regra de inferência. Se a regra de inferência for modus ponens,
procedemos como na demonstração de \olref[ded]{thm:deduction-thm}. Se a
regra de inferência for \QR, sabemos que $!B \ident !C \lif \lforall[x][!D(x)]$
e que !!a{formula} da forma $!C \lif !D(a)$ aparece anteriormente na
!!{derivation}, em que $a$ não ocorre em~$!C$, $!A$ nem em $\Gamma$. Temos,
portanto,
\begin{align*}
  \Gamma \cup \{!A\} & \Proves !C \lif !D(a),\\
  \intertext{e a hipótese de indução se aplica, isto é, temos}
    \Gamma & \Proves !A \lif (!C \lif !D(a)).\\
  \intertext{Por}
  & \Proves (!A \lif (!C \lif !D(a))) \lif ((!A \land !C) \lif !D(a))\\
  \intertext{e por modus ponens obtemos}
  \Gamma & \Proves (!A \land !C) \lif !D(a).\\
  \intertext{Como a condição de variável própria continua satisfeita, podemos acrescentar a essa !!{derivation} um passo justificado por \QR e obter}
    \Gamma & \Proves (!A \land !C) \lif \lforall[x][!D(x)].\\
    \intertext{Também temos}
    & \Proves ((!A \land !C) \lif \lforall[x][!D(x)]) \lif (!A \lif (!C \lif \lforall[x][!D(x)]),\\
    \intertext{portanto, por modus ponens,}
    \Gamma & \Proves !A \lif (!C \lif \lforall[x][!D(x)]),
\end{align*}
isto é, $\Gamma \Proves !A \lif !B$.

Deixamos como exercício o caso em que $!B$ é justificada pela regra \QR,
mas tem a forma $\lexists[x][!D(x)] \lif !C$.
\end{proof}

\begin{prob}
  Complete a demonstração de \olref[fol][axd][ddq]{thm:deduction-thm-q}.
\end{prob}

\end{document}
